%% Define repeatable tests, conditions, equipment, and pass/fail criteria.
\chapter{Validation methods}
\label{sec:testing}

Select validation methods appropriate to the implemented functionality,
available equipment, and project requirements. Describe tests performed on
the final, completed prototype. For each test, describe its purpose, setup,
procedure, data recorded, and objective pass/fail criterion. Where relevant,
document the equipment, test conditions, repeated measurements, and
measurement uncertainty. Provide sufficient detail for the test to be
reproduced. Report deviations from the planned procedure.

% Describe tests of sensor performance against references or specifications.
\section{Sensor-characteristic tests}

Describe how relevant sensor characteristics or performance were evaluated,
such as measurement accuracy, repeatability, response time, detection
reliability, or recognition performance. Explain the test conditions,
methods, and reference measurements where applicable. Distinguish measured
performance from manufacturer specifications and acknowledge characteristics
that could not be verified.

% Describe tests of system functions, required conditions, and error handling.
\section{Demonstrator-function tests}

List the procedures used to verify acquisition, processing, display or other
outputs, feedback, error handling, and operation across the required
conditions. Where applicable, evaluate startup behavior, repeated operation,
and stability over time to identify intermittent failures or inconsistent
behavior. Link each procedure to requirement IDs in
Table~\ref{tab:requirements}.

% Explain how user understanding is evaluated and how responses are recorded.
\section{User-understanding evaluation (optional)}

If a user-understanding evaluation was conducted, describe the participants
in aggregate without identifying them, the tasks or questions given, the test
environment, how responses were recorded, and the criteria for judging
understanding. Address any relevant consent, privacy, or ethical considerations.